\section{实验结果与分析}

\subsection{运行结果}

三种算法在 pcb442 实例上的 10 次运行结果如下。

\begin{table}[H]
\centering
\caption{10 次运行最优解对比}
\begin{tabular}{crrr}
\toprule
运行 & dEA & HdEA & HdEA-MP \\
\midrule
1 & 51690 & 57166 & 73349 \\
2 & 51693 & 57062 & 81099 \\
3 & 51549 & 57866 & 74888 \\
4 & 51394 & 56273 & 71605 \\
5 & 51724 & 55544 & 80696 \\
6 & 51571 & 55593 & 75038 \\
7 & 51554 & 57080 & 76812 \\
8 & 51696 & 56932 & 77613 \\
9 & 51488 & 56261 & 73665 \\
10 & 51702 & 56182 & 75596 \\
\midrule
\textbf{均值} & \textbf{51606} & \textbf{56596} & \textbf{76036} \\
\textbf{标准差} & \textbf{112} & \textbf{746} & \textbf{3081} \\
\textbf{最小值} & \textbf{51394} & \textbf{55544} & \textbf{71605} \\
\textbf{最大值} & \textbf{51724} & \textbf{57866} & \textbf{81099} \\
\bottomrule
\end{tabular}
\end{table}

\subsection{统计显著性分析}

对三种算法的 10 次结果两两进行配对 t 检验（$\alpha=0.05$），结果见表 \ref{tab:ttest}。

\begin{table}[H]
\centering
\caption{配对 t 检验结果}
\label{tab:ttest}
\begin{tabular}{lccc}
\toprule
对比 & $t$ 值 & $p$ 值 & 是否显著 \\
\midrule
dEA vs HdEA & -20.93 & $<0.001$ & 显著，dEA 更优 \\
dEA vs HdEA-MP & -25.71 & $<0.001$ & 显著，dEA 更优 \\
HdEA vs HdEA-MP & -19.14 & $<0.001$ & 显著，HdEA 更优 \\
\bottomrule
\end{tabular}
\end{table}

所有对比的 $p$ 值均远小于 0.05，说明三种算法的性能差异具有统计显著性。

\subsection{可视化对比}

图 \ref{fig:boxplot} 的箱线图和图 \ref{fig:barplot} 的柱状图直观展示了三种算法的分布特征。

\begin{figure}[H]
\centering
\includegraphics[width=0.85\textwidth]{boxplot.png}
\caption{箱线图：三种算法 10 次运行结果分布}
\label{fig:boxplot}
\end{figure}

\begin{figure}[H]
\centering
\includegraphics[width=0.85\textwidth]{barplot.png}
\caption{柱状图：三种算法均值和标准差对比（误差线表示 $\pm 1$ 标准差）}
\label{fig:barplot}
\end{figure}

\subsection{结果分析}

从实验结果可以看出，三种并行进化算法在求解质量上呈现出明显而一致的梯度：dEA 最优，HdEA 次之，HdEA-MP 最差，且两两之间的差异均通过了配对 t 检验的显著性验证（$p < 0.001$）。

dEA 以均值 51606、标准差仅 112 的成绩在求解质量上显著领先于另外两种算法。其核心优势在于全局共享种群的设计——200 个个体在同一个进化池中进行竞争、选择和交配，提供了最为充分的基因多样性，从而使种群能够持续探索搜索空间，避免陷入局部最优。从标准差来看，dEA 在 10 次运行中表现极为稳定，波动幅度（51394--51724）不超过 330，远小于 HdEA 和 HdEA-MP。

HdEA 的求解质量（均值 56596，标准差 746）明显低于 dEA，尽管其采用了更复杂的层级架构。导致这一结果的根本原因是岛屿规模过小：每个子种群仅有 50 个个体，在 442 城市的庞大搜索空间中，这一规模难以维持足够的遗传多样性。各岛屿在进化早期便快速收敛至不同的局部最优，虽然每 20 代进行一次精英迁移（将各岛最优个体汇集到 Master 再广播），但迁移仅限于最优个体的交换，对于恢复种群整体多样性的作用十分有限，最终整体最优解仍然受限于各岛屿的局部收敛质量。

HdEA-MP 的表现进一步印证了上述分析。在 HdEA 基础上增加环形迁移机制后，解质量反而出现了明显下降（均值 76036，标准差 3081）。值得注意的是，HdEA-MP 的标准差高达 3081，是三种算法中方差最大的，说明环形迁移在某些运行中确实引入了额外的基因多样性——部分运行的最优解可达到 71605。然而，由于每岛种群规模过小，迁移过来的优质基因无法在种群内有效传播和继承，最终进化结果呈现出高度的不确定性。即便是 HdEA-MP 所有运行中的最优结果（71605），也仍逊于 dEA 最差的结果（51724）。

在精度与时间这一经典权衡维度上，三种算法呈现出不同的取舍模式。dEA 每代需要两次全局集体通信（Bcast + Gatherv），通信开销约为 0.22 ms/代；HdEA 和 HdEA-MP 的大部分世代在本地独立进化，仅每 20 代进行一次通信，折合每代通信开销仅约 0.07 ms。从运行时间来看，dEA 的总耗时（约 44 s）约为 HdEA-MP（约 13 s）的 3 倍。然而 dEA 的路径距离相比 HdEA-MP 降低了约 32\%，这一解质量的提升远超额外通信所付出的时间成本。在精度优先的实际应用场景中，dEA 无疑是更为合理的选择。
